% A conditional evaluated *inside another conditional's operand*.
%
% tex.web 498 keeps open conditionals on a stack (cond_ptr), and that is not
% bookkeeping for its own sake:
%
%   \if e\if Exe\else x\fi T\else F\fi
%
% The inner \if runs while the outer is scanning its second operand. It takes
% its false branch and leaves its \fi pending in the stream. When the outer
% then skips its own false branch, that pending \fi is the first one it meets
% -- and it is not the outer's. TeX pops the inner conditional and keeps
% skipping; a plain depth counter stops there, ends the outer early, and runs
% its true branch instead.
%
% expl3 nests conditionals inside conditionals everywhere. l3fp's exponent
% test is exactly this shape:
%
%   \if:w e\if:w E\exp_not:N #1e\else: \exp_not:N #1\fi: ...
%
% so every floating-point literal was read as though it carried an exponent.
%
% The \noexpand cases matter separately: \exp_not:N sits in front of nearly
% every operand expl3 tests, and it is a *non*-conditional primitive expanded
% mid-scan -- it must leave the enclosing conditional's depth untouched.
\input plain
\catcode`\:=11 \catcode`\_=11
\def\M#1{\message{<#1>}}
% A conditional evaluated while an enclosing one scans its operands. The
% inner leaves its \fi pending, and that \fi is the first the outer meets
% when it skips its own false branch -- it is not the outer's.
\M{a=\if e\if Exe\else x\fi T\else F\fi}
\M{b=\if e\if EEe\else x\fi T\else F\fi}
\M{c=\if x\if Exe\else x\fi T\else F\fi}
% Two levels of nesting in operand position.
\M{d=\if e\if E\if AAx\else e\fi e\else x\fi T\else F\fi}
% The same through \noexpand, as expl3 writes it.
\let\if:w\if \let\exp_not:N\noexpand \let\fi:\fi \let\else:\else
\let\s__mark\relax
\def\test#1{\if:w e\if:w E\exp_not:N #1e\else: \exp_not:N #1\fi: TRUE\else: FALSE\fi:}
\M{e=\test\s__mark}
\M{f=\test e}
\M{g=\test E}
\M{h=\test x}
\def\mac{q}
\M{i=\test\mac}
% Repeated, to catch depth drifting across evaluations.
\M{j=\test\s__mark}
\M{k=\test E}
\M{l=\test\s__mark}
% Nested in an \ifnum operand.
\M{m=\ifnum\ifnum1>0 1\else 2\fi=1 T\else F\fi}
\M{n=\ifnum\ifnum1>2 1\else 2\fi=1 T\else F\fi}
% \ifcase with a conditional choosing the case.
\M{o=\ifcase\ifnum1>0 1\else 0\fi zero\or one\else other\fi}
\bye
